\documentclass[12pt]{article}
\usepackage[utf8]{inputenc}
\usepackage[T1]{fontenc}
\usepackage[margin=1in]{geometry}
\usepackage{setspace}
\usepackage{xcolor}
\usepackage[colorlinks=true, urlcolor=blue!60!black, linkcolor=blue!60!black]{hyperref}
\setlength{\parindent}{0em}
\setlength{\parskip}{0.6em}
\renewcommand{\baselinestretch}{1.0}
\begin{document}
\begin{center}
{\LARGE \textit{Seila Law LLC v. Consumer Financial Protection Bureau}, 591 U.S. 197 (2020)\par}
\vspace{0.5em}
{\large Case Brief --- Two-Page Edition\par}
\end{center}

\section{Memory Jogger}
A single agency head wielding significant executive power cannot be shielded from at-will presidential removal; the CFPB's removal restriction was unconstitutional but severable (\href{https://www.law.cornell.edu/supremecourt/text/19-7}{Cornell LII}).

\section{Detailed Case Facts}
Dodd-Frank created the CFPB with broad rulemaking, enforcement, and adjudicatory power, headed by a single Director with a five-year term, removable only for ``inefficiency, neglect of duty, or malfeasance in office,'' and funded through the Federal Reserve. Seila Law, a debt-relief law firm served with a civil investigative demand, challenged the structure.

\section{Procedural History}
The district court enforced the demand and the Ninth Circuit affirmed under \textit{Humphrey's Executor} and \textit{Morrison}. The Government agreed the restriction was invalid, so the Court appointed an amicus to defend the judgment.

\section{Judicial Votes}
On removal, 5--4: Roberts for Thomas, Alito, Gorsuch, and Kavanaugh; Kagan, with Ginsburg, Breyer, and Sotomayor, dissented. On severability, 7--2: Roberts's plurality plus Kagan's four; Thomas and Gorsuch dissented.

\section{Holding}
The President's executive power generally includes at-will removal, subject only to the \textit{Humphrey's Executor} exception for multimember expert bodies and the \textit{Morrison} exception for limited inferior officers. The CFPB Director's protection is unconstitutional but severable; the case was remanded on ratification.

\section{Analysis of Opinions}
\textbf{Majority:} The CFPB fits neither exception; a single-Director independent agency lacks historical footing and defies a design that concentrates accountability in one elected President, a problem aggravated by the long term and independent funding. \textbf{Plurality:} the express severability clause keeps the rest of the Act operative. \textbf{Thomas:} overrule \textit{Humphrey's Executor}, and do not ``excise'' statutes. \textbf{Kagan:} Congress may structure offices under the Necessary and Proper Clause; history is full of insulated financial regulators; and a single Director is easier to control than a board.

\section{Examples: Future Applications}
\textbf{Same-Side:} the single FHFA Director's protection falls (\textit{Collins v. Yellen}). \textbf{Opposite-Side:} a bipartisan five-member adjudicatory commission survives under \textit{Humphrey's Executor}. \textbf{Fence-Sitter:} the modern FTC, a multimember body wielding substantial executive power, later addressed in \textit{Trump v. Slaughter} (2026).

\section{Critique}
Unitary-executive originalists say Roberts stopped halfway; revisionist historians reply that early Congresses routinely regulated tenure. Functionalists call the single-head versus multimember line arbitrary and warn of threats to insulated financial regulators.

\section{Key Quotations}
\begin{quote}
``Under our Constitution, the `executive Power'---all of it---is `vested in a President,' who must `take Care that the Laws be faithfully executed.' Art. II, \S 1, cl. 1; id., \S 3.''
\end{quote}
\noindent --- Chief Justice Roberts, for the Court.

\end{document}
